\subsection{Modified phase behavior}
Boussinesq approximation approximates phase density with constant densities but
buoyancy terms are left alone.
The density difference is essential in gravitational driven melt segaration\cite{MagmaDynamicsEnthalpyMethod}.
This time we keep the difference in liquid and solid density, stating 
$\rho_s \neq \rho_f$.
We still nondimensionalize in the similar way:
\begin{equation}
\begin{aligned}
\mathcal{C} &= \phi_2 + \frac{\rho_f}{\rho_s}\phi_f \chi_f \\
\mathcal{H} &= (\phi_1 + \phi_2)\mathcal{T} + \phi_f\mathcal{T} + \phi_f L_f \\
            &= (1-\phi_f)\mathcal{T} + 
               \frac{\rho_f}{\rho_s}\phi_f(\frac{L}{c_p} + \mathcal{T})\\
            &= [1-(1-\frac{\rho_f}{\rho_s})\phi_f]\mathcal{T} + 
               \frac{\rho_f}{\rho_s}\phi_f\mathcal{L}
\end{aligned}
\end{equation}
We then separate phase regions similarly.

\subsubsection{Region I}
Single phase solidus region: olivine.
This condition was not affected since no liquid presents.

\subsubsection{Region II}
Two-phase sub-solidus: olivine + opx.
This region was not affected by differentiate liquid and solid densities
since no liquid presents in this region.

\subsubsection{Region III}
This is three phase eutectic region, all three phases are present.
In this phase region, temperature is fixed at the eutectic temperature,
where the dimensionless temperature is 0.
\begin{equation}
		  \phi_f = \frac{\rho_s}{\rho_f}\frac{\mathcal{H}}{L} \quad \phi_2 = \mathcal{C} - \frac{\rho_f}{\rho_s}\phi_f\chi_f(\mathcal{T}, \mathcal{P}) = \mathcal{C} - \frac{\mathcal{H}}{L}  \quad \phi_1 = 1-\phi_f - \phi_2 \quad \text{and} \quad \mathcal{T} = \mathcal{T}_e = \gamma P
\end{equation}
Where $\chi_f(\mathcal{T}, \mathcal{P}) = \chi_e(\mathcal{P})$, all melting is 
eutectic melting. Mass fraction is fixed as eutectic mass fraction.

\subsubsection{Region IV}
In this super-eutectic two-phase region, opx has been exhausted in eutectic melting.
The volumetric fractions can be determined as:
\begin{equation}
\phi_2 = 0 \quad \phi_f = \frac{\rho_s}{\rho_f}\frac{\mathcal{C}}{\chi_f(\mathcal{T}, \mathcal{P})}
\end{equation}
The original quadratic equation has now been altered as:
\begin{equation}
F(\mathcal{T}) = \mathcal{H} - [1-(1-\frac{\rho_f}{\rho_s})\frac{\rho_s}{\rho_f}\frac{\mathcal{C}}{\chi_f}]\mathcal{T} - \frac{\mathcal{C}}{\chi_f(\mathcal{T}, \mathcal{P})} L = 0 
\end{equation}
We use simple linear relationship between mass fraction and temperature as 
$\chi_f = \chi_e(\mathcal{T}_m-\mathcal{T})$.
\begin{equation}
\mathcal{T}^2 - (\mathcal{H} + \mathcal{T}_m - 
(1-\frac{\rho_f}{\rho_s})\frac{\rho_s}{\rho_f}\frac{\mathcal{C}}{\chi_e})\mathcal{T}
+\mathcal{T}_m \mathcal{H} - \frac{\mathcal{C} L}{\chi_e} = 0
\end{equation}

\subsubsection{Region VI}
In this single phase  super-liquidus region, we can get volume fractions and temperature as usual:
\begin{equation}
		  \phi_1 = \phi_2 = 0 \quad \phi_f = 1 \quad \text{and} \quad \mathcal{T} = \frac{\rho_s}{\rho_f}\mathcal{H} - L
\end{equation}

\subsubsection{Two lines}
Similarly, we can calculate two lines separating regions III, IV and VI using relationship between composition and bulk 
mass fraction.
\begin{equation}
		  \mathcal{C} = \frac{\rho_f}{\rho_s} \chi
\end{equation}
Hence we can substitute $\chi$ with $\mathcal{C}$ simply.
